\section{Conclusion}
\label{sec:conclusion}

Where a market supplies a monthly forward strip and no options, the parts of
an option value that are identified and the parts that are assumed can be kept
apart by construction. An equality-constrained hourly curve carries the
identified level, a regime-switching residual with time-varying transitions
carries the distributional assumptions, and a centring term ensures that the
second cannot alter the first under any parameter vector or change of measure.

Three implications follow for valuation in the Turkish market and in markets
like it. The forward level of the unquoted front month bounds the precision of
any short-dated option value, and a value is more informative when reported
with the range that this level induces. Residual parameters inherited from
another estimation need to be tested on a series that matches their use; here
they were individually wrong but correct in the stationary variance that
governs the reported maturities, and that agreement held from 2023 to 2025 and
failed in 2026, when dispersion stopped scaling with the price level. Finally,
forward-anchored valuation inherits the predictive accuracy of the strip, which
was low and unstable in sign over 2022--2026, and the strip moved too rarely to
hedge with.

A residual whose dispersion adapts to the price regime, and a premium on
regime risk calibrated once option-like prices become observable, are the
extensions that the evidence points to.
